\chapter{Conclusion}\label{chapter:conclusion}

  \section{Key Findings and Results}
    This thesis provides multiple results and findings relevant for the research area of \ac{EEG} based classification, especially in the context of \ac{MGR} and hybrid acquisition methods. For these purposes, a novel \ac{EEG} dataset tailored for \ac{MGR} is introduced, being the first of its kind. The balanced dataset consisting of 100 songs per genre with ten genres in total is available in standard file formats and is easily reproducable. Furthermore, a hybrid \ac{EEG} headset was utilised to measure the signal, fostering novel approaches that could make their way into everyday life. As the headset also acquires other biosignals than the neural activity, those were also measured and are accessible in each file.\par
    Going hand in hand with the \ac{EEG} dataset, a custom denoising and preprocessing pipeline is presented that works on fewer electrodes and unvonventional placement of those. By using an \ac{ICA} based approach combined with a score based rejection mechanism, it is possible to detect artifacts in the measured signal to eliminate them from the data. \par
    Apart from the creation of a custom \ac{EEG} datasets, \ac{MGR} has been performed not only on those data but also on the well-known \ac{DEAP} dataset to have another source of brain signals as comparision. Overall, it is possible to achieve accuracy and F1-scores better than chance level on classification for each subject independently. This infers that there are neural patterns specific to musical genres that can be captured by \ac{ML} methods, being a promising finding for further research. Moreover, a fluctuating trend and variance in performance metrics is observable, leading to the conclusion that there are inconsitencies in music genre patterns and data quality across sessions for the \ac{DEAP} and the headset data respectively. Nonetheless, it is significant that the majority of optimal models are using a sigmoid kernel, implying that this function is able to capture most of the genre specific patterns. \par
  \section{Problems and Limitations}
    Although this thesis provides key findings and conclusion for the area of \ac{EEG} based \ac{MGR}, there are still some problems and limitations that isolate out of this thesis. The first one lies in the hybrid \ac{EEG} headset approach, as it still remains questionable whether the unconventional placing can measure the neural activity well. The reason for this concern is that state-of-the-art \ac{EEG} devices place the electrodes on top of the scalp with a large amount (mainly 32-64) of electrodes. The headset, however, has not only fewer electrodes (seven) but they are also not place on the scalp but around the ear. \par
    Another limitation arises from the fact that each model is optimized and trained for each subject individually. This strategy may work for tasks that require dependency on the subject (e.\,g. personalised therapy), but does not give away information about a possible general purpose model being able to perform well on unseen subjects. \par
    Despite the fact that the performances are above chance level, the overall results are still only slightly better than this. Especially mult-class classification for the headset dataset remains a difficult task as the differentiation between more than two genres seems hard to capture. \par

  \section{Future Work}
    To overcome the mentioned limitations and problems, further research in the area of \ac{EEG} based \ac{MGR} has to be conducted. First, to verify the quality of the \ac{EEG} signal recorded by the headset, it would be vital to record the neural activity with a state-of-the-art \ac{EEG} device and standardised denoising steps on the same songs and session setting. This way, it would be possible to compare both signals with each other to make statements about the signal's quality of the hybrid headset. Moreover, it would be beneficial to further extend the headset dataset to have more versatile data. Besides adding more songs, different subjects could be recorded. For example, it would be interesting to see how the signal differ between musicians and non-musicians.\par
    Other future work consists of different a classification evaluation by not only optimising per subject/session but for all data collectively. This way, more evidence can be given how well \ac{ML} methods can generalise on different subjects and/or different environments. Combined with that, it would be interesting to see how well different \ac{ML} techniques perform on this task. Especially Deep Learning architectures have the potential to yield even better results as they are widely used in other applications regarding \ac{EEG}. \par
    Last but not least, since the headset also measures other biosignals, those can be included in the pipeline. It can not only help capturing artifacts better but also offer valuable information that can improve performance.
